% Chapter 1: The Objects of Learning
% Rewritten against the design-lab Lean 4 kernel (FLT_Proofs.*):
% typed definitions, full proofs isomorphic to the verified arguments,
% audit-driven corrections, and a closing set of open problems.
% concept_class = Profile G (hub node, parallax),
% inductive_bias = Profile D (negative result centerpiece),
% MDL/MML/Bayesian = Profile E (analogy triangle obstruction analysis)

\chapter{The Objects of Learning}
\label{ch:objects}

Every learning problem begins with the same question: given data drawn from an unknown source, find a rule that predicts well on future data. This chapter introduces the mathematical vocabulary for making that question precise.

The vocabulary divides into three groups of unequal importance. The first group, domain, label, concept, hypothesis, is atomic: these are the sets and functions from which everything else is built. They require only brief definitions. The second group, concept class, hypothesis space, version space, inductive bias, carries genuine mathematical content, because the relationships between these objects determine what is learnable; we prove the load-bearing facts in full. The third group, Bayesian inference, minimum description length, minimum message length, appears to be vocabulary but is actually a network of near-miss analogies, each encoding a different answer to the same question about model selection.

Every typed object in this chapter has a machine-checked counterpart, named where one exists.

\section{Atomic Vocabulary}
\label{sec:atoms}

The following objects are the type-level substrate of learning theory. They are presented together because none of them, individually, contains any mathematical surprise.

\begin{defn}[Domain, Label, Concept]
\label{def:domain-label-concept}
A \emph{domain} $X$ is a set whose elements are called \emph{instances}. A \emph{label set} $Y$ is a set of possible outputs; in binary classification, $Y = \{0, 1\}$. A \emph{concept} is a function $c : X \to Y$.\formal{FLT_Proofs.Basic.Concept}

Equivalently, when $Y = \{0,1\}$, a concept $c$ can be identified with the subset $\{x \in X : c(x) = 1\} \subseteq X$. Both perspectives, function and set, are used throughout the literature. We adopt the function view as primary, matching the typed definition $\leanname{Concept X Y} := X \to Y$, and switch to the set view when it simplifies combinatorial arguments (as in shattering, \cref{def:shatter-vc}).
\end{defn}

\begin{defn}[Hypothesis, Target Concept, Proper/Improper]
\label{def:hypothesis-target}
A \emph{hypothesis} $h : X \to Y$ is a candidate prediction rule that a learning algorithm might output; as a type it is a concept by another name, $\leanname{Hypothesis X Y} := \leanname{Concept X Y}$, and the distinction is semantic: a hypothesis is the learner's guess, a concept is ground truth. The \emph{target concept} $c^\ast \in \mathcal{C}$ is the specific concept that generated the training data, an element of the subtype $\{c : c \in \mathcal{C}\}$. Learning is \emph{proper} when the learner's hypothesis space equals the concept class, $\mathcal{H} = \mathcal{C}$, and \emph{improper} when it may use a larger space $\mathcal{H} \supseteq \mathcal{C}$.\formal{FLT_Proofs.Basic.IsProper}
\end{defn}

\begin{defn}[Realizable and Agnostic]
\label{def:realizable-objects}
The learning problem $(\mathcal{C}, \mathcal{H})$ is \emph{realizable} when every target is available to the learner, $\mathcal{C} \subseteq \mathcal{H}$.\formal{FLT_Proofs.Basic.Realizable} It is \emph{agnostic} when no such containment is assumed; the agnostic setting is precisely the \emph{absence} of a realizability hypothesis, and the learner competes against the best hypothesis in $\mathcal{H}$, whatever its error, since a perfect one may be unavailable.
\end{defn}

An \emph{alphabet} $\Sigma$ is a finite symbol set; it appears only in the formal language chapters (\Cref{ch:automata,ch:gold}) and plays no role in statistical learning theory.

\medskip

These definitions are unremarkable individually. The mathematical content begins when we ask how they compose.

% ================================================================
% CONCEPT CLASS: Hub Node, Profile G (Framework Bridge)
% Centerpiece: the node everything measures
% ================================================================

\section{The Concept Class: Where All Paradigms Converge}
\label{sec:concept-class}

\begin{defn}[Concept Class]
\label{def:concept-class}
A \emph{concept class} $\mathcal{C} \subseteq Y^X$ is a collection of concepts over a common domain; as a type, $\leanname{ConceptClass X Y} := \mathrm{Set}\,(\leanname{Concept X Y})$.\formal{FLT_Proofs.Basic.ConceptClass} The class is the primary object of study: every major question in learning theory is a question about concept classes.
\end{defn}

This definition is simple. Its importance is not. The concept class is the single object on which the chapter's deepest theorem turns: a family of apparently different demands one might place on it (learnability from few examples, finite shattering, bounded compression, vanishing correlation with noise, sub-exponential growth, finite sample size) are secretly one condition, and the class is the vertex at which they coincide (Figure~\ref{fig:concept-class-hub}).\formal{FLT_Proofs.Theorem.PAC.fundamental_theorem}

\begin{figure}[ht]
\centering
\begin{tikzpicture}[
    pivot/.style={draw, very thick, rounded corners, fill=defblue!20, font=\footnotesize\bfseries, minimum width=2.1cm, minimum height=0.95cm, align=center},
    face/.style={draw, thick, rounded corners, fill=defblue!8, font=\scriptsize, minimum width=1.85cm, minimum height=0.74cm, align=center},
    onlinenode/.style={draw, thick, rounded corners, fill=onlineblue!15, font=\scriptsize, minimum width=1.45cm, minimum height=0.62cm, align=center},
    frontier/.style={draw, dashed, rounded corners, fill=edgeorange!7, font=\scriptsize, minimum width=1.6cm, minimum height=0.58cm, align=center},
    chipgood/.style={draw, circle, thick, fill=sepgreen!25, font=\tiny, inner sep=0.4pt, minimum size=0.56cm},
    chipbad/.style={draw, circle, thick, fill=thmred!18, font=\tiny, inner sep=0.4pt, minimum size=0.56cm},
    equiv/.style={{Stealth[length=1.5mm]}-{Stealth[length=1.5mm]}, semithick, defblue},
    strictbridge/.style={-{Stealth[length=1.8mm]}, thick, onlineblue!75!black},
    fedge/.style={-{Stealth[length=1.3mm]}, dashed, edgeorange!85!black},
    elab/.style={font=\tiny, text=notegray, inner sep=1.1pt, fill=white},
    flab/.style={font=\tiny\itshape, text=edgeorange!65!black, inner sep=1pt}
]
% measurability gate
\draw[dashed, thick, notegray] (0,-0.1) ellipse (3.15cm and 2.75cm);
\node[notegray, font=\tiny] at (0,-2.35) {measurability gate};
% pivot: the loop-closing measure on the carrier
\node[pivot] (vc) at (0,0) {$\mathcal{C}$\\$\VC(\mathcal{C})<\infty$};
% verified equivalence faces (conjuncts of the Fundamental Theorem)
\node[face] (pac)  at (-2.05,-1.5) {PAC\\learnable};
\node[face] (rad)  at (-2.5, 0.95) {Rademacher\\$\to 0$};
\node[face] (grow) at (0, 2.2)     {sub-exp.\\growth};
\node[face] (comp) at (2.5, 0.95)  {compression\\(with info)};
\node[face] (samp) at (2.05,-1.5)  {finite sample\\complexity};
% two-way machine-checked equivalences
\draw[equiv] (vc) -- node[elab]{Fund.\ Thm} (pac);
\draw[equiv] (vc) -- (rad);
\draw[equiv] (vc) -- node[elab]{Sauer--Shelah} (grow);
\draw[equiv] (vc) -- node[elab]{Moran--Yehud.} (comp);
\draw[equiv] (vc) -- node[elab]{$\Theta(\VC)$} (samp);
\draw[equiv] (pac) -- (rad);
% nested verified loop: online learning
\node[onlinenode] (onl) at (-2.55,-3.05) {online\\learnable};
\node[onlinenode] (ld)  at (-0.65,-3.05) {Littlestone\\dim};
\draw[equiv, onlineblue!70!black] (onl) -- (ld);
\draw[strictbridge] (onl) to[bend left=14] node[elab, pos=0.58]{strict $\not\Leftarrow$} (pac);
% attention chips straddling the gate
\node[chipgood] (soft) at (1.95,1.6) {soft};
\node[chipbad]  (hard) at (3.0,2.45) {hard};
\node[font=\tiny, text=sepgreen!55!black, anchor=east] at (1.58,1.6) {softmax};
\node[font=\tiny, text=thmred!70!black, anchor=west] at (3.34,2.45) {argmax};
% generality frontier: literature / open / known-to-break
\node[frontier] (ds)    at (0, 3.35)    {multiclass / DS};
\node[frontier] (noinf) at (4.05, 0.95) {compression,\\no side info};
\node[frontier] (agn)   at (-4.5,-1.5)  {agnostic PAC};
\node[frontier, draw=thmred!70!black, fill=thmred!5] (part) at (-4.5, 0.5) {partial\\concepts};
\draw[fedge] (ds)    -- node[flab,right]{lit.\ '22} (grow);
\draw[fedge] (noinf) -- node[flab,above]{open} (comp);
\draw[fedge] (agn)   -- node[flab,below]{lit.} (pac);
\draw[fedge, thmred!70!black] (part) -- node[flab,above,text=thmred!75!black]{VC$\,\not\Leftrightarrow\,$PAC} (pac);
% legend
\node[font=\tiny, anchor=west, text=notegray] at (-4.5,-3.7) {%
  \textcolor{defblue}{\rule[0.45ex]{0.5cm}{0.8pt}}~verified equivalence \quad
  \textcolor{edgeorange!85!black}{\rule[0.45ex]{0.15cm}{0.8pt}\,\rule[0.45ex]{0.15cm}{0.8pt}}~literature / open \quad
  \textcolor{onlineblue!75!black}{\rule[0.45ex]{0.5cm}{0.8pt}}~strict one-way};
\end{tikzpicture}
\caption[The concept class as the vertex of an equivalence]{\textbf{The concept class as the vertex of an equivalence.} Several demands on a class of concepts (PAC learnability, finite shattering, bounded compression, vanishing Rademacher correlation, sub-exponential growth, finite sample complexity) are one property wearing several faces, and the class is the object on which they coincide. A tighter wheel for online learning sits inside, joined by a one-way bridge since the online demand implies the statistical one but not conversely; the whole rests on a measurability gate, and a dashed frontier marks coincidences expected but not yet verified, with the edge for partial concepts known to break.}
\label{fig:concept-class-hub}
\end{figure}

Each measure asks a different question about the same object:

\begin{itemize}[leftmargin=*,itemsep=2pt]
\item \textbf{VC dimension} asks: how many points can $\mathcal{C}$ shatter? (Determines PAC learnability.)
\item \textbf{Littlestone dimension} asks: how deep a mistake tree can $\mathcal{C}$ support? (Determines online learnability.)
\item \textbf{DS dimension} asks: what pseudo-cubes exist in $\mathcal{C}$'s one-inclusion hypergraph? (Determines multiclass learnability.)
\item \textbf{Teaching dimension} asks: how many examples suffice for a teacher to identify each $c \in \mathcal{C}$? (Measures cooperative communication.)
\item \textbf{Rademacher complexity} asks: how well can $\mathcal{C}$ correlate with random noise? (Controls generalization bounds.)
\end{itemize}

These 13 measures are not independent: many are related by \rel{upper\_bounds}, \rel{lower\_bounds}, and \rel{characterizes} edges in the graph. But none is reducible to another; each captures an aspect of $\mathcal{C}$'s complexity the others miss. That a single set of functions supports 13 inequivalent complexity measures is one of the structural surprises of the field.

A single theorem is what installs the concept class as the hub, in the place a paradigm or a measure might naively be expected to occupy. To state it we need the first of the measures.

\begin{defn}[Shattering and VC dimension]
\label{def:shatter-vc}
A finite set $S \subseteq X$ is \emph{shattered} by $\mathcal{C}$ if every labeling of $S$ is realized by some concept in $\mathcal{C}$:
\[
\text{$S$ shattered} \iff \forall\, f : S \to \{0,1\},\ \exists\, c \in \mathcal{C},\ \forall\, x \in S,\ c(x) = f(x).
\]\formal{FLT_Proofs.Complexity.VCDimension.Shatters}
The \emph{VC dimension} $\VC(\mathcal{C}) \in \N \cup \{\infty\}$ is the supremum of the cardinalities of shattered sets:
\[
\VC(\mathcal{C}) = \sup\{\,|S| : S \subseteq X \text{ finite and shattered by } \mathcal{C}\,\}.
\]\formal{FLT_Proofs.Complexity.VCDimension.VCDim}
\end{defn}

The full theory of VC dimension is the subject of \Cref{ch:dimensions}; here we need only the definition and one elementary closure property, used repeatedly below: a subset of a shattered set is shattered. If $T \subseteq S$ and $S$ is shattered, then any labeling of $T$ extends to a labeling of $S$, which some concept realizes; restricting to $T$ proves $T$ shattered.\formal{FLT_Proofs.Complexity.Littlestone.Shatters.subset}

\begin{thm}[VC characterization of PAC learnability]
\label{thm:vc-char}
For a concept class $\mathcal{C}$ over a measurable domain $X$ (subject to the mild measurability conditions of \Cref{ch:generalization}),
\[
\mathcal{C} \text{ is PAC learnable} \iff \VC(\mathcal{C}) < \infty.
\]\formal{FLT_Proofs.Theorem.PAC.vc_characterization}
\end{thm}

The proof occupies \Cref{ch:pac} (sufficiency, via empirical risk minimization and uniform convergence) and \Cref{ch:dimensions} (necessity, via a hard distribution on a shattered set); it is the reason the concept class is the hub of \cref{fig:concept-class-hub}. The class, through a single combinatorial invariant of itself, decides its own learnability. The equivalence extends to a five-way characterization: PAC learnability, finite VC dimension, the existence of a bounded sample compression scheme, uniform convergence of Rademacher complexity, and a finite sample-complexity bound are all equivalent.\formal{FLT_Proofs.Theorem.PAC.fundamental_theorem} We meet the compression and Rademacher faces in \Cref{ch:compression,ch:generalization}.

\begin{remark}[The equivalence is asymmetric]
\label{rmk:vc-asymmetry}
\Cref{thm:vc-char} is an ``iff'', but its two halves are proved by opposite kinds of argument. The sufficiency direction ($\VC < \infty \Rightarrow$ learnable) is constructive: it exhibits an algorithm, empirical risk minimization, and bounds its error. The necessity direction (learnable $\Rightarrow \VC < \infty$) is non-constructive: from an infinite VC dimension it produces a distribution on which every learner fails, without exhibiting any particular hard instance the learner must confront in advance. The two directions cross the boundary between combinatorics and measure theory in opposite directions, and whether the necessity direction admits a constructive proof is open (\cref{op:asymmetry}).
\end{remark}

We can now compute the invariant on the smallest interesting classes, and in doing so witness a separation that recurs throughout the book.

\begin{example}[Four concept classes of increasing complexity]
\label{ex:concept-classes}
Fix $X = \R$.
\begin{enumerate}[label=(\alph*),leftmargin=*,itemsep=3pt]
\item \textbf{Thresholds:} $\mathcal{C}_1 = \{x \mapsto \ind[x \geq \theta] : \theta \in \R\}$. Here $\VC(\mathcal{C}_1) = 1$ and $\Ldim(\mathcal{C}_1) = \infty$ (\cref{prop:threshold-vc,prop:threshold-ldim}).
\item \textbf{Intervals:} $\mathcal{C}_2 = \{x \mapsto \ind[a \leq x \leq b] : a \leq b \in \R\}$. Here $\VC(\mathcal{C}_2) = 2$: the two-point set $\{0,1\}$ is shattered (the four labelings are realized by intervals disjoint from both, containing only $0$, only $1$, or both), while no three points $x_1 < x_2 < x_3$ are, because the labeling $(1,0,1)$ would require an interval containing $x_1$ and $x_3$ but not $x_2$.
\item \textbf{Unions of $k$ intervals:} $\mathcal{C}_3^{(k)}$, with $\VC(\mathcal{C}_3^{(k)}) = 2k$ by the same argument applied $k$ times; the detailed counting is deferred to \Cref{ch:dimensions}.
\item \textbf{All measurable functions:} $\mathcal{C}_4 = \{0,1\}^X$, with $\VC(\mathcal{C}_4) = \infty$. By \cref{thm:vc-char} it is not PAC learnable; by the same token it is not online learnable nor identifiable in the limit. The No-Free-Lunch theorem (\cref{thm:nfl}) applies to it in full force.
\end{enumerate}
The thresholds deserve attention: finite VC dimension but infinite Littlestone dimension. This single class witnesses the separation between PAC and online learning (\Cref{ch:separations}). We prove both values now, because the proofs are the prototypes for the shattering and adversary arguments used throughout.
\end{example}

\begin{prop}[Thresholds have VC dimension one]
\label{prop:threshold-vc}
The class $\mathcal{C}_1 = \{x \mapsto \ind[x \geq \theta] : \theta \in \R\}$ has $\VC(\mathcal{C}_1) = 1$.
\end{prop}

\begin{proof}
\emph{At least one.} The singleton $\{x_0\}$ is shattered: the labeling $x_0 \mapsto 1$ is realized by any $\theta \leq x_0$, and $x_0 \mapsto 0$ by any $\theta > x_0$.

\emph{Not two.} Let $x_1 < x_2$. The labeling $f(x_1) = 1$, $f(x_2) = 0$ asks for a threshold $\theta$ with $x_1 \geq \theta$ and $x_2 < \theta$, hence $x_2 < \theta \leq x_1$, contradicting $x_1 < x_2$. So no two-point set is shattered. By \cref{def:shatter-vc}, $\VC(\mathcal{C}_1) = 1$.
\end{proof}

\begin{prop}[Thresholds have infinite Littlestone dimension]
\label{prop:threshold-ldim}
The class $\mathcal{C}_1$ admits, against any online learner, a realizable sequence forcing an unbounded number of mistakes; hence $\Ldim(\mathcal{C}_1) = \infty$.
\end{prop}

\begin{proof}
We exhibit an adversary that forces a mistake every round while keeping the transcript consistent with some threshold. The adversary maintains an open interval $(a,b)$ of still-admissible thresholds, initialized to $(0,1)$. At each round it presents the midpoint $x = (a+b)/2$ and reads the learner's prediction $\hat{y} = \ind[x \geq \hat\theta]$. It then declares the true label to be $1 - \hat{y}$ and updates:
\begin{itemize}[nosep]
\item if $\hat{y} = 1$, declare label $0$ (meaning $x < \theta$, i.e.\ $\theta > x$) and set $(a,b) \leftarrow (x, b)$;
\item if $\hat{y} = 0$, declare label $1$ (meaning $x \geq \theta$, i.e.\ $\theta \leq x$) and set $(a,b) \leftarrow (a, x)$.
\end{itemize}
After the update the interval $(a,b)$ remains nonempty (the midpoint strictly separates it), and \emph{every} $\theta$ in it is consistent with the entire transcript so far: each declared label $\ind[x_i \geq \theta]$ is correct for such $\theta$ by construction. Thus the sequence is realizable by the class at every finite stage, yet the learner errs at every round. Since the number of rounds is unbounded, the realizable mistake bound is infinite, and therefore $\Ldim(\mathcal{C}_1) = \infty$.\formal{FLT_Proofs.Theorem.Separation.adversary_from_shatters}
\end{proof}

\Cref{prop:threshold-vc,prop:threshold-ldim} together discharge the claim of \cref{ex:concept-classes}(a): the same class is trivially PAC learnable (one example pins the threshold up to $\varepsilon$) yet admits no finite online mistake bound. Because \cref{prop:threshold-ldim} forces unbounded mistakes against every learner, the gap is a genuine separation; we return to it in \Cref{ch:separations}.

% ================================================================
% HYPOTHESIS SPACE: Profile B-light
% Centerpiece: proper vs improper tension
% ================================================================

\section{Hypothesis Spaces and the Proper--Improper Divide}
\label{sec:hypothesis-space}

\begin{defn}[Hypothesis Space]
\label{def:hypothesis-space}
A \emph{hypothesis space} $\mathcal{H} \subseteq Y^X$ is the set of functions available to the learner as candidate outputs. As a type it is structurally identical to a concept class, $\leanname{HypothesisSpace X Y} := \mathrm{Set}\,(\leanname{Concept X Y})$, and semantically distinct: $\mathcal{C}$ is the ground-truth collection, $\mathcal{H}$ is what the learner may return. When $\mathcal{H} = \mathcal{C}$, learning is \emph{proper}; when $\mathcal{H} \supsetneq \mathcal{C}$, learning is \emph{improper}.
\end{defn}

The proper--improper distinction carries real weight. Its computational consequences are among the sharpest in the theory:

\begin{enumerate}[leftmargin=*,itemsep=3pt]
\item \textbf{Improper learning can be exponentially cheaper.} There exist concept classes where proper learning requires exponentially more samples or computation than improper learning~\cite{PittValiant1988}. The graph records this as a \rel{does\_not\_imply} edge from \nde{pac\_learning} to \nde{exact\_learning} (\Cref{ch:separations}).

\item \textbf{Computational hardness often depends on properness.} Cryptographic hardness results (Kearns--Valiant, 1994) typically establish hardness of \emph{proper} learning; improper learners may circumvent the barrier by using richer representations~\cite{KearnsValiant1994}.

\item \textbf{The characterization is representation-independent.} \Cref{thm:vc-char} holds regardless of whether learning is proper or improper. What decides learnability is $\VC(\mathcal{C})$, the VC dimension of the \emph{target} class; properties of the learner's hypothesis space $\mathcal{H}$ leave the verdict untouched. The learner's freedom to range outside $\mathcal{C}$ moves the cost while leaving possibility fixed.
\end{enumerate}

The simplest object built from $\mathcal{H}$ and data together is the set of survivors.

\begin{defn}[Version Space]
\label{def:version-space}
Given a hypothesis space $\mathcal{H}$ and a sample $S = \{(x_i, y_i)\}_{i=1}^m$, the \emph{version space} is
\[
\mathrm{VS}_{\mathcal{H},S} = \{h \in \mathcal{H} : h(x_i) = y_i \text{ for all } i\},
\]
the set of hypotheses consistent with all observed data.\formal{FLT_Proofs.Complexity.GameInfra.versionSpace} The verified development carries a bundled $\leanname{VersionSpace}$ structure packaging $\mathcal{H}$, the data, and the consistency obligations.\formal{FLT_Proofs.Basic.VersionSpace}
\end{defn}

Version-space elimination (returning any element of $\mathrm{VS}_{\mathcal{H},S}$) is among the most natural learning algorithms. Its basic behavior is captured by four propositions, each a one-line consequence of the definition and each a prototype for the convergence arguments of \Cref{ch:gold}.

\begin{prop}[Soundness, persistence, and monotonicity of the version space]
\label{prop:version-space}
Fix $\mathcal{H}$, a target $c \in \mathcal{H}$, and a sample $S$ labeled by $c$. Then
\begin{enumerate}[label=(\roman*),leftmargin=*,itemsep=2pt]
\item $\mathrm{VS}_{\mathcal{H},S} \subseteq \mathcal{H}$;\formal{FLT_Proofs.Complexity.GameInfra.versionSpace_subset}
\item $c \in \mathrm{VS}_{\mathcal{H},S}$;\formal{FLT_Proofs.Complexity.GameInfra.target_in_versionSpace}
\item $\mathrm{VS}_{\mathcal{H},\,S \cup \{(x,y)\}} \subseteq \mathrm{VS}_{\mathcal{H},S}$ for any further labeled example $(x,y)$;\formal{FLT_Proofs.Complexity.GameInfra.versionSpace_append}
\item $\Ldim(\mathrm{VS}_{\mathcal{H},S}) \leq \Ldim(\mathcal{H})$.\formal{FLT_Proofs.Complexity.GameInfra.ldim_versionSpace_le}
\end{enumerate}
\end{prop}

\begin{proof}
(i) An element of $\mathrm{VS}_{\mathcal{H},S}$ is by definition a member of $\mathcal{H}$ satisfying the consistency constraints; forgetting the constraints gives membership in $\mathcal{H}$.

(ii) The target $c$ lies in $\mathcal{H}$ by hypothesis, and since $S$ is labeled by $c$ we have $c(x_i) = y_i$ for every $i$; so $c$ meets both defining conditions and lies in the version space. In particular the version space is never empty in the realizable case.

(iii) Adding a labeled example $(x,y)$ adjoins the constraint $h(x) = y$ to the defining conjunction. Any $h$ satisfying the larger conjunction satisfies the smaller one, so the consistent set can only shrink.

(iv) By (i), $\mathrm{VS}_{\mathcal{H},S} \subseteq \mathcal{H}$. The Littlestone dimension is monotone under inclusion: a mistake tree shattered by a subclass is shattered, unchanged, by the superclass (each labeling realized inside the subclass is realized inside the superclass). Taking the supremum of shattered-tree depths over the smaller class gives a value no larger than over the larger one.
\end{proof}

Item (iii) is the engine of identification in the limit: as data accumulates, the version space contracts monotonically, and the question of \Cref{ch:gold} is whether it contracts to the target. Item (iv) says that this contraction never makes the residual problem harder in the online sense, a fact the Standard Optimal Algorithm exploits to achieve the optimal mistake bound.

The trouble with version-space elimination is computational. For most hypothesis classes of interest, the version space cannot be efficiently represented, enumerated, or sampled from. Two dual invariants quantify aspects of this difficulty. The \emph{eluder dimension}~\cite{RussoVanRoy2013} measures how long an adversary can keep producing examples on which consistent hypotheses still disagree, controlling exploration--exploitation tradeoffs. The \emph{star number} measures a complementary, purely combinatorial quantity.

\begin{defn}[Star number]
\label{def:star-number}
The \emph{star number} of $\mathcal{C}$ is the largest $d$ for which there is a set $S = \{x_1,\dots,x_d\}$ such that each point can be isolated: for every $x \in S$ there is a concept $c_x \in \mathcal{C}$ with $c_x(x) = 1$ and $c_x(y) = 0$ for all other $y \in S$.\formal{FLT_Proofs.Complexity.VCDimension.StarNumber}
\end{defn}

The star number is never smaller than the VC dimension, and the proof is immediate from shattering.

\begin{prop}[Star number dominates VC dimension]
\label{prop:star-ge-vc}
For every concept class, $\VC(\mathcal{C}) \leq \mathrm{star}(\mathcal{C})$.
\end{prop}

\begin{proof}
Let $S$ be a shattered set of size $d = \VC(\mathcal{C})$ (if $\VC(\mathcal{C}) = \infty$, apply the argument to shattered sets of every size). For each $x \in S$ the single labeling ``$x \mapsto 1$, all other points of $S$ $\mapsto 0$'' is one of the labelings realized by shattering, so there is a concept isolating $x$ within $S$. Hence $S$ witnesses the star condition, and $\mathrm{star}(\mathcal{C}) \geq |S| = d$.
\end{proof}

\begin{warning}[The reverse inequality fails, and so does the folklore equivalence]
\label{rmk:star-not-equiv}
It is tempting to believe that finiteness of the star number is \emph{equivalent} to finiteness of the VC dimension. It is not. Consider the class of singletons on an infinite domain, $\mathcal{C} = \{\ind[\cdot = a] : a \in X\}$. Its VC dimension is $1$: a single point $\{a\}$ is shattered ($\ind[\cdot=a]$ labels it $1$, any other singleton labels it $0$), but no pair $\{a,b\}$ is, because labeling both points $1$ would require a single concept true at $a$ and $b$, which no singleton is. Yet its star number is infinite: for \emph{any} finite set $S$ and any $x \in S$, the concept $\ind[\cdot = x]$ isolates $x$ within $S$, so every finite $S$ is a star set. Thus $\VC = 1$ while $\mathrm{star} = \infty$. Finiteness of the star number is therefore strictly stronger than finiteness of the VC dimension; the implication runs one way only (\cref{prop:star-ge-vc}). The star number governs the label complexity of \emph{active} learning, where this gap is exactly the point~\cite{HannekeYang2015}; we return to it in \Cref{ch:dimensions}. Identifying which finite-VC classes have finite star number is open (\cref{op:star}).
\end{warning}

% ================================================================
% INDUCTIVE BIAS: Profile D (Negative Result)
% Centerpiece: NFL theorem makes this NECESSARY
% ================================================================

\section{Inductive Bias and the No-Free-Lunch Barrier}
\label{sec:inductive-bias}

\begin{defn}[Inductive Bias]
\label{def:inductive-bias}
The \emph{inductive bias} of a learner is the set of assumptions, explicit or implicit, that it uses to generalize from training data to unseen instances.
\end{defn}

This informal definition is the one usually given, and it suggests that inductive bias is not a single mathematical object. That suggestion is half right. The \emph{content} of a bias does change across paradigms, as the table below records. But the \emph{form} is uniform, and can be typed: a bias is a hypothesis space together with a real-valued preference ranking on concepts, structured so that nothing outside the space is preferred to anything inside it.

\begin{defn}[Inductive bias, typed]
\label{def:inductive-bias-typed}
An \emph{inductive bias} over $X, Y$ is a triple $(\mathcal{H}, \mathrm{pref}, \pi)$ where $\mathcal{H} \subseteq Y^X$ is a hypothesis space, $\mathrm{pref} : Y^X \to \R$ is a preference score (lower is more preferred), and $\pi$ is the constraint that every concept outside $\mathcal{H}$ is scored at least as high as every concept inside it: $h \notin \mathcal{H} \Rightarrow \mathrm{pref}(h') \leq \mathrm{pref}(h)$ for all $h' \in \mathcal{H}$.\formal{FLT_Proofs.Basic.InductiveBias}
\end{defn}

This single structure subsumes the standard readings. Reading $\mathrm{pref}(h) = -\log P(h)$ recovers the Bayesian prior; reading $\mathrm{pref}(h)$ as description length recovers minimum description length; reading it as the index of the first class in a nested hierarchy containing $h$ recovers structural risk minimization. The differences between these readings are real (\cref{sec:model-selection}), but they are differences of \emph{which} preference score, inside one type.

\begin{center}\small
\begin{tabularx}{\textwidth}{@{}l l X@{}}
\toprule
\textbf{Paradigm} & \textbf{Inductive bias formalized as} & \textbf{Graph edge} \\
\midrule
PAC & Restriction of $\mathcal{H}$ to finite VC dimension & \rel{used\_in\_proof} by \nde{nfl\_theorem} \\
Bayesian & Prior $P(h)$ over hypotheses & \rel{analogy} from \nde{bayesian\_inference} \\
SRM & Nested class hierarchy $\mathcal{H}_1 \subset \mathcal{H}_2 \subset \cdots$ & \rel{analogy} from \nde{srm} \\
Computational & Cryptographic assumption restricting efficient search & \rel{requires\_assumption} by \nde{computational\_hardness} \\
\bottomrule
\end{tabularx}
\end{center}

Inductive bias is a \emph{mathematical necessity}, forced on every learner, and the reason is the No-Free-Lunch theorem. We prove it through its combinatorial engine, then read off the classical statement and the necessity of bias as corollaries.

\begin{lemma}[Shattering forces a hard concept]
\label{lem:nfl-engine}
Let $T \subseteq X$ be shattered by $\mathcal{C}$ with $|T| > 2m$. For any $m$ instances $x_1,\dots,x_m$ and any hypothesis $h : X \to \{0,1\}$, there is a concept $c \in \mathcal{C}$ such that
\begin{enumerate}[label=(\roman*),leftmargin=*,nosep]
\item $c$ agrees with $h$ on every sample point that lies in $T$: $c(x_i) = h(x_i)$ whenever $x_i \in T$; and
\item $c$ disagrees with $h$ on more than a quarter of $T$: $\bigl|\{x \in T : c(x) \neq h(x)\}\bigr| > |T|/4$.
\end{enumerate}\formal{FLT_Proofs.Complexity.Generalization.nfl_per_sample_shattered}
\end{lemma}

\begin{proof}
Define an adversarial labeling $f : T \to \{0,1\}$ that copies $h$ on the seen points and flips it on the rest:
\[
f(x) = \begin{cases} h(x) & \text{if } x \in \{x_1,\dots,x_m\},\\ 1 - h(x) & \text{otherwise.}\end{cases}
\]
Because $T$ is shattered, some $c \in \mathcal{C}$ realizes $f$ on $T$, i.e.\ $c(x) = f(x)$ for all $x \in T$. On a seen point $x_i \in T$ we have $c(x_i) = f(x_i) = h(x_i)$, giving (i). On an unseen point $x \in T$ we have $c(x) = f(x) = 1 - h(x) \neq h(x)$, so $c$ disagrees with $h$ on the entire unseen part of $T$. The unseen part has size at least $|T| - m$, since the $m$ sample points can cover at most $m$ elements of $T$. Finally, $|T| > 2m$ gives $|T| - m > |T|/2 > |T|/4$, so the disagreement count exceeds $|T|/4$, which is (ii).
\end{proof}

\begin{thm}[No Free Lunch {\cite{Wolpert1996}}]
\label{thm:nfl}
Let $A$ be any learning algorithm that observes $m$ labeled examples, and suppose $\mathcal{C}$ shatters some set $T$ with $|T| \geq 2m$. Then there is a target concept $c^\ast \in \mathcal{C}$ and the uniform distribution $D$ on $T$ such that
\[
\E_{S \sim D^m}\bigl[\risk_D(A(S))\bigr] \geq \tfrac{1}{4},
\]
where $S$ is labeled by $c^\ast$. In particular, no learner does better than this on its worst target: averaging cannot rescue what shattering has spoiled.
\end{thm}

\begin{proof}
Take $|T| = 2m$ and let $D$ be uniform on $T$. Shattering realizes all $2^{|T|}$ labelings of $T$ inside $\mathcal{C}$; let $c$ range uniformly over these $2^{2m}$ concepts. It suffices to bound the average risk over $c$ from below by $1/4$, since the maximum over targets is at least the average. Writing the risk as an average of disagreement indicators over $T$ and exchanging the (finite) expectations,
\[
\E_{c}\,\E_{S \sim D^m}\bigl[\risk_D(A(S))\bigr]
= \frac{1}{|T|}\sum_{x \in T} \E_{S}\,\E_{c}\bigl[\ind[A(S)(x) \neq c(x)]\bigr].
\]
Fix a draw $S$ of $m$ points and a point $x \in T$ not among them. The learner's output $A(S)$ depends on $c$ only through the labels at the seen points, so $A(S)(x)$ is independent of $c(x)$; and $c(x)$ is a uniform bit. Hence the inner expectation is exactly $\tfrac{1}{2}$ for every unseen $x$. The unseen points number at least $|T| - m = m = |T|/2$, so the average is at least $\tfrac{1}{|T|}\cdot \tfrac{|T|}{2}\cdot\tfrac{1}{2} = \tfrac14$. Some target $c^\ast$ therefore attains $\E_S[\risk_D(A(S))] \geq \tfrac14$. (The per-sample form of this argument, which builds the hard target adversarially in place of averaging over targets, is \cref{lem:nfl-engine}.)
\end{proof}

\begin{cor}[Inductive bias is mandatory]
\label{cor:bias-mandatory}
A concept class of infinite VC dimension is not PAC learnable; in particular the class $\{0,1\}^X$ of all binary functions on an infinite domain is not learnable. Consequently every learnable class is a proper subset of $\{0,1\}^X$, and the choice of which subset is exactly the inductive bias.
\end{cor}

\begin{proof}
If $\VC(\mathcal{C}) = \infty$, then $\mathcal{C}$ shatters sets of every size, so for each sample budget $m$ it shatters some $T$ with $|T| \geq 2m$. \Cref{thm:nfl} then supplies, for every learner, a target and distribution on which the expected error is at least $1/4$; taking $\varepsilon, \delta < 1/4$ shows the PAC criterion cannot be met. Hence infinite VC dimension precludes PAC learning, the necessity half of \cref{thm:vc-char}, here proved directly.\formal{FLT_Proofs.Complexity.Generalization.vcdim_infinite_not_pac} Since $\{0,1\}^X$ has infinite VC dimension, it is unlearnable, so any learnable class omits some functions; the omission \emph{is} the assumption that those functions are not the target.
\end{proof}

\begin{remark}[No-Free-Lunch as a boundary of askability]
\label{rmk:nfl-pl}
\Cref{cor:bias-mandatory} is more than a negative result; it relocates inductive bias from the engineering of a learner to the precondition of the question. Before a class is fixed, ``which functions are efficiently learnable?'' has no nontrivial answer, because none are. Inductive bias is what makes the learnability question askable at all. The graph encodes this as a \rel{requires\_assumption} dependency from \nde{inductive\_bias} into \nde{computational\_hardness}: the hardness theory presupposes the restriction that No-Free-Lunch shows to be unavoidable.
\end{remark}

% ================================================================
% MODEL SELECTION: Profile E (Analogy Triangle)
% Centerpiece: MDL/MML/Bayesian near-miss network
% ================================================================

\section{Three Faces of Model Selection}
\label{sec:model-selection}

Bayesian inference, minimum description length, and minimum message length are often described informally as ``the same idea.'' They are not. Each formalizes a different answer to the question ``which hypothesis should I prefer?''; the differences between them are more instructive than any individual definition. Each is, in the sense of \cref{def:inductive-bias-typed}, a choice of preference score; the analogies among them are analogies among scores, and the obstructions are precisely where the scores fail to agree.

\subsection{The Definitions}

\begin{defn}[Bayesian Inference]
\label{def:bayesian}
Given a prior $P(h)$ over $\mathcal{H}$ and a likelihood $P(S \mid h)$, the posterior is $P(h \mid S) \propto P(S \mid h) \cdot P(h)$. The Bayesian learner outputs $\hat{h} = \arg\max_{h} P(h \mid S)$ (MAP) or the full posterior $P(\cdot \mid S)$. As a type, a Bayesian learner bundles a prior with the posterior-update rule and converts to an ordinary batch learner by MAP read-off.\formal{FLT_Proofs.Learner.Bayesian.BayesianLearner}
\end{defn}

\begin{defn}[Minimum Description Length {\cite{Rissanen1978}}]
\label{def:mdl}
Select $h$ minimizing $L(h) + L(S \mid h)$, where $L(\cdot)$ denotes description length under a fixed prefix-free code. The first term penalizes complexity; the second penalizes misfit.
\end{defn}

\begin{defn}[Minimum Message Length {\cite{WallaceBoulton1968}}]
\label{def:mml}
Select $h$ minimizing the total message length $\ell(h) + \ell(S \mid h)$ required to transmit the hypothesis and then the data, under a two-part coding scheme. Wallace--Freeman (1987): $\ell(h) = -\log P(h) + \frac{1}{2}\log |F(h)| + \text{const}$, where $F(h)$ is the Fisher information.
\end{defn}

\subsection{The Analogy Network}

These three methods are usually treated as one idea. The deeper structure, made precise by the chapter's compression characterization, is that they are approximations to a single machine-checked equivalence: a class is learnable exactly when its data compress to a bounded summary that reconstructs every consistent labelling (Figure~\ref{fig:model-selection-triangle}).\formal{FLT_Proofs.Theorem.PAC.fundamental_vc_compression} The pairwise analogies among the methods remain genuine but unproven, each blocked by a different obstruction.

\begin{figure}[ht]
\centering
\begin{tikzpicture}[
    spine/.style={draw, very thick, rounded corners, fill=defblue!14, font=\small, minimum width=2.3cm, minimum height=0.85cm, align=center},
    sub/.style={draw, thick, rounded corners, fill=defblue!5, font=\scriptsize, minimum width=2.1cm, minimum height=0.7cm, align=center},
    bridge/.style={draw, thick, rounded corners, fill=sepgreen!12, font=\scriptsize, minimum width=2.0cm, minimum height=0.7cm, align=center},
    method/.style={draw, thick, rounded corners, fill=analogypurple!10, font=\small\bfseries, minimum width=1.7cm, minimum height=0.7cm, align=center},
    vedge/.style={-{Stealth[length=1.7mm]}, semithick, defblue},
    centerpiece/.style={{Stealth[length=2.2mm]}-{Stealth[length=2.2mm]}, line width=1.3pt, defblue},
    gedge/.style={-{Stealth[length=1.7mm]}, semithick, sepgreen!70!black},
    dedge/.style={-{Stealth[length=1.5mm]}, dashed, edgeorange!85!black},
    obs/.style={dashed, analogypurple!55, semithick},
    openedge/.style={-{Stealth[length=1.5mm]}, dashed, thmred!75!black},
    elab/.style={font=\tiny, text=notegray, inner sep=1.3pt, fill=white},
    olab/.style={font=\tiny, text=analogypurple!70!black, inner sep=1.5pt, fill=white, align=center}
]
% spine: the machine-checked equivalence
\node[spine] (vc)   at (-2.3,0) {finite VC\\dimension};
\node[spine] (comp) at ( 2.3,0) {bounded compression\\(with side info)};
\draw[centerpiece] (vc) -- node[elab,above=1pt]{learnable $\Leftrightarrow$ compressible} node[elab,below=1pt]{Moran--Yehudayoff} (comp);
% side-information axis: proven vs open asymmetry
\node[sub] (noinf) at (-2.3,-2.0) {compression\\(no side info)};
\draw[vedge]    ([xshift=-4pt]noinf.north) -- node[elab,left]{VC $\le 2(k{+}1)^2$} ([xshift=-4pt]vc.south);
\draw[openedge] ([xshift=4pt]vc.south) -- node[elab,right]{$O(d)$? open} ([xshift=4pt]noinf.north);
% PAC-Bayes bridge (the one proven cross-edge)
\node[bridge] (pb) at (0,1.7) {prior-weighted\\complexity};
\draw[gedge] (pb) -- node[elab,right=1pt]{PAC-Bayes} (vc);
% model-selection principles as approximating instances
\node[method] (bayes) at (-3.3,3.5) {Bayesian};
\node[method] (mdl)   at ( 0,  3.5) {MDL};
\node[method] (mml)   at ( 3.3,3.5) {MML};
\draw[vedge] (bayes) -- node[elab,left=1pt]{MAP $\to$ batch} (pb);
\draw[dedge] (mdl) -- node[elab,right]{description length} (comp);
\draw[dedge] (mml) to[bend left=8] node[elab,right]{two-part code} (comp);
% obstruction arcs: no equivalence theorem exists
\draw[obs] (bayes) -- node[olab]{type\\mismatch} (mdl);
\draw[obs] (mdl) -- node[olab]{missing\\witness} (mml);
\draw[obs] (bayes) to[bend left=18] node[olab]{missing equivalence witness} (mml);
\node[font=\scriptsize\itshape, text=notegray] at (0,5.0) {model-selection principles (each a choice of inductive bias)};
% legend
\node[font=\tiny, anchor=west, text=notegray] at (-4.6,-2.95) {%
  \textcolor{defblue}{\rule[0.45ex]{0.5cm}{1pt}}~machine-checked \quad
  \textcolor{edgeorange!85!black}{\rule[0.45ex]{0.15cm}{0.8pt}\,\rule[0.45ex]{0.15cm}{0.8pt}}~literature \quad
  \textcolor{thmred!75!black}{\rule[0.45ex]{0.15cm}{0.8pt}\,\rule[0.45ex]{0.15cm}{0.8pt}}~open};
\end{tikzpicture}
\caption[The compression core of model selection]{\textbf{The compression core of model selection.} Bayesian inference, minimum description length, and minimum message length are usually treated as one idea. The figure separates them around the fact the chapter proves: a class is learnable exactly when its data admit a bounded summary reconstructing every consistent labelling. That equivalence is the spine; the three principles hang off it as approximations, with every link among them unproven. Whether the auxiliary tag can be removed while keeping the summary short is open (\cref{op:compression-objects}).}
\label{fig:model-selection-triangle}
\end{figure}

\subsection{Why They Are Not the Same}

\begin{obstbox}
\textbf{Bayesian $\leftrightarrow$ MML:} obstruction type \emph{missing\_equivalence\_witness}. MML's two-part code corresponds to Bayesian MAP estimation when the Fisher information correction term is constant. But for models with varying curvature in parameter space, MML departs from Bayesian posteriors. No theorem establishes their equivalence beyond the constant-Fisher case.
\end{obstbox}

\begin{obstbox}
\textbf{MML $\leftrightarrow$ MDL:} obstruction type \emph{missing\_equivalence\_witness}. Both minimize a two-part code length. But MML uses Shannon-optimal coding relative to a prior (subjective), while MDL uses a universal code based on Kolmogorov complexity (objective). In the limit of infinite data, both converge to the true model under standard regularity conditions, but their finite-sample behavior differs.
\end{obstbox}

\begin{obstbox}
\textbf{MDL $\leftrightarrow$ Bayesian:} obstruction type \emph{type\_mismatch}. MDL uses description lengths (integers under a code); Bayesian inference uses probability densities (reals under a measure). The connection $L(h) = -\log P(h)$ identifies description lengths with log-probabilities, but this identification breaks down when the code is not prefix-free or the prior is improper.
\end{obstbox}

That three different formalizations of ``prefer the simpler hypothesis'' are \emph{not} equivalent, and that the reasons for non-equivalence are themselves different (missing witness, missing witness, and type mismatch), is a structural feature of learning theory that no single definition can convey. It is visible only in the graph.

One bridge across this network is an actual theorem. The preference score of a Bayesian or MDL learner controls generalization through the PAC-Bayes bound: for any prior fixed before seeing the data and any posterior, the expected risk is bounded by the empirical risk plus a complexity term governed by the $\KL$ divergence from prior to posterior~\cite{McAllester1999}, proved here for finite hypothesis spaces.\formal{FLT_Proofs.Theorem.PACBayes.pac_bayes_finite} The obstruction boxes concern the three \emph{procedures}, which are not interchangeable; PAC-Bayes concerns the \emph{quantity} they share, a prior-weighted complexity, and shows it bounds error.

\subsection{Compression, Occam, and the Fundamental Theorem}

The MDL idea (prefer the hypothesis that compresses the data) connects directly to \cref{thm:vc-char}. A \emph{sample compression scheme} of size $k$ reconstructs any consistent labeling from a subsequence of at most $k$ of its own examples. Compression and finite VC dimension are equivalent.

\begin{thm}[Compression characterizes finite VC dimension {\cite{MoranYehudayoff2016}}]
\label{thm:compression-char}
A concept class has finite VC dimension if and only if it admits a sample compression scheme of bounded size (with side information):
\[
\VC(\mathcal{C}) < \infty \iff \exists\, k,\ \mathcal{C} \text{ has a size-}k\text{ compression scheme with side information.}
\]\formal{FLT_Proofs.Theorem.PAC.fundamental_vc_compression}
\end{thm}

We prove this in \Cref{ch:compression}; here it closes the circle. Through \cref{thm:vc-char,thm:compression-char}, ``prefer the hypothesis that compresses'' and ``prefer the learnable hypothesis'' are the same condition on the class. The qualification ``with side information'' carries weight: whether every class of VC dimension $d$ admits a compression scheme of size $O(d)$ \emph{without} side information is the Littlestone--Warmuth conjecture, still open (\cref{op:compression-objects}).\formal{FLT_Proofs.Complexity.Structures.Open_NoInfoCompressionStrengthening}

\subsection{Structural Risk Minimization}

\begin{defn}[Structural Risk Minimization {\cite{Vapnik1998}}]
\label{def:srm}
Given a nested sequence of hypothesis classes $\mathcal{H}_1 \subset \mathcal{H}_2 \subset \cdots$ with $\VC(\mathcal{H}_k) = d_k$, select $k^\ast$ minimizing $\emprisk_{\mathcal{H}_k}(h) + \mathrm{penalty}(d_k, m, \delta)$.
\end{defn}

SRM is the operational realization of inductive bias within VC theory: instead of choosing a single $\mathcal{H}$, choose a hierarchy indexed by complexity, and let the data select the level. In the language of \cref{def:inductive-bias-typed}, SRM's preference score is the index $k$ of the first class containing $h$; the nesting guarantees the monotonicity condition $\pi$. The graph records an \rel{analogy} edge from \nde{srm} to \nde{inductive\_bias} with a \emph{type\_mismatch} obstruction, because SRM constructs a bias by structural nesting while ``inductive bias'' as a concept is broader than any one construction.

\begin{defn}[Algorithmic Probability {\cite{Solomonoff1964}}]
\label{def:alg-prob}
The \emph{Solomonoff prior} assigns to each string $x$ the probability $M(x) = \sum_{p : U(p)=x} 2^{-|p|}$, summing over all programs $p$ that produce $x$ on a universal Turing machine $U$. This is the ``prior that assigns higher probability to simpler explanations'' in the most literal possible sense.
\end{defn}

Algorithmic probability connects to Kolmogorov complexity through the coding theorem, $K(x) = -\log M(x) + O(1)$, and provides the information-theoretic foundation for MDL. Its detailed treatment belongs to \Cref{ch:compression}, where it interacts with compression schemes and Occam's razor; the explicit constant in the coding theorem is itself a target for formalization (\cref{op:solomonoff}).

\section{What This Chapter Established}
\label{sec:ch1-summary}

The base vocabulary of learning theory is small: domain, label, concept, concept class, hypothesis space, version space. Each is a type, and each typed object has a machine-checked counterpart. At this level alone, four structural features are already visible.

\begin{enumerate}[leftmargin=*,itemsep=3pt]
\item The concept class is a hub node: the same object is measured by 13 inequivalent complexity measures, and one of them, the VC dimension, decides learnability outright (\cref{thm:vc-char}). The multiplicity earns its keep: each of the 13 reflects a genuinely different aspect of a single object.

\item Inductive bias is not optional. The No-Free-Lunch theorem (\cref{thm:nfl}), proved here through its shattering engine, converts bias from a design choice into a precondition for the learnability question to be askable at all (\cref{cor:bias-mandatory}).

\item Negative results are content, and rigor catches folklore. The asserted equivalence between finiteness of the star number and of the VC dimension is false, with singletons as an explicit witness (\cref{rmk:star-not-equiv}); the true relationship is a one-sided inequality (\cref{prop:star-ge-vc}).

\item The three standard model-selection principles, Bayesian, MDL, and MML, form a triangle of near-miss analogies whose obstructions differ pair by pair, while a genuine theorem (PAC-Bayes) and a genuine equivalence (compression $\leftrightarrow$ finite VC, \cref{thm:compression-char}) sit nearby and are often confused with them.
\end{enumerate}

These patterns recur throughout the book: hub nodes, negative results as content, analogy obstructions, and the discipline of separating what is proved from what is merely said.

% ================================================================
% TRANSFORMER CONNECTION (retrofit): attention as a concept class,
% well-behavedness, and the soft/hard = Borel/analytic boundary.
% ================================================================

\section{The Objects of a Transformer}
\label{sec:ch1-transformer}

The vocabulary of this chapter is not confined to the classical examples. A single softmax-attention head, read as a binary predictor, is a concept in the sense of \cref{def:domain-label-concept}, a function $X \to \{0,1\}$, and letting its parameters vary produces a concept class in the sense of \cref{def:concept-class}. The question returns with a sharp answer for the central object of modern machine learning: \emph{is the attention class a legitimate concept class, one to which the VC characterization applies?} Answering it requires exactly the measurability that \cref{thm:vc-char} deferred to \Cref{ch:generalization}; this section makes those ``mild measurability conditions'' concrete and shows them non-vacuous.

\begin{defn}[The softmax-attention concept class]
\label{def:attention-concept-class}
Fix an embedding dimension $d$ and a context of $n_K$ key--value pairs.  Take the domain to be the query space $X = \R^d$.  For keys $K = (k_1,\dots,k_{n_K})$ and values $V = (v_1,\dots,v_{n_K})$, the head maps a query $x$ to the readout $r_{K,V}(x) = \sum_{i=1}^{n_K} \mathrm{softmax}(\langle x, k_i\rangle)_i\, v_i$, and thresholding at zero gives a concept
\[
  c_{K,V}(x) \;=\; \ind[\,r_{K,V}(x) > 0\,] \;:\; \R^d \to \{0,1\}.
\]
The \emph{softmax-attention concept class} is $\mathcal{A}_{d,n_K} = \{\, c_{K,V} : K, V \,\}$, the range of the parametrized family.\formal{TLT.Bridge.softAttentionConcept}
\end{defn}

So far this is only a particular concept class among many. Its first non-trivial property is that it is a \emph{well-behaved} object: the measurability hypothesis that \cref{thm:vc-char} assumes is satisfied.

\begin{prop}[Attention is a well-behaved object]
\label{prop:attention-wellbehaved}
For every $d$ and $n_K$, the class $\mathcal{A}_{d,n_K}$ is a well-behaved VC-measurability target: the bad events over which the uniform-convergence argument of \Cref{ch:generalization} quantifies are measurable, so $\mathcal{A}_{d,n_K}$ is a bona-fide concept class for the VC characterization.\formal{TLT.Bridge.softAttention_wellBehaved}
\end{prop}

\begin{proof}
A parametrized concept family is a well-behaved VC-measurability target as soon as the evaluation map $(K,V,x) \mapsto c_{K,V}(x)$ is jointly measurable.  Here that map factors through the readout $(K,V,x) \mapsto r_{K,V}(x)$, which is a finite sum of products of softmax weights and values; each softmax weight is a continuous function of its arguments, so the readout is jointly continuous and therefore Borel measurable.  The concept is the indicator of the open condition $\{r_{K,V}(x) > 0\}$, whose preimage is the measurable set $\{(K,V,x) : 0 < r_{K,V}(x)\}$.  Joint measurability follows, and with it well-behavedness.  The base case is the same statement for the simplest concept class of all: the singleton-indicator class on a measurable set is well-behaved,\formal{TLT.Tame.singletonClassOn_wellBehavedVCMeasTarget} and attention's well-behavedness is the parametrized lift of that fact.
\end{proof}

The proof used continuity of the softmax in an essential way.  Far from a cosmetic change, removing it carries the attention class clear across a boundary in descriptive set theory.

\begin{prop}[Softmax is the boundary of measurability]
\label{prop:soft-hard-borel}
Assume the existence of an analytic set of reals that is not Borel (a standard fact of classical descriptive set theory).  Then softmax attention is well-behaved at every $(d,n_K)$, \emph{and} there is a hard-attention witness, a continuous score map and two measurable experts, whose argmax-routed head has a bad event that is not measurable.\formal{TLT.Bridge.soft_vs_hard_attention_measurabilitySeparation}
\end{prop}

\begin{proof}[Proof idea]
The first conjunct is \cref{prop:attention-wellbehaved}.  For the second, replace the softmax by a hard $\arg\max$ over the scores.  The set of queries routed to a given key is then a level set of an $\arg\max$, and one can arrange, using a continuous reduction from an analytic non-Borel set, that the resulting bad event is the continuous preimage of that analytic set, hence analytic but not Borel, hence not measurable.  The two routes to attention thus sit on opposite sides of the Borel/analytic line: softmax stays measurable; $\arg\max$ can fall out of the measurable world entirely.
\end{proof}

\begin{remark}[Where the deferred conditions bite]
\label{rmk:attention-measurability}
\Cref{prop:soft-hard-borel} is the concrete content behind the parenthetical of \cref{thm:vc-char}.  The ``mild measurability conditions'' are not automatically met by every natural hypothesis class: hard attention violates them, and the violation is genuine: an analytic-not-Borel bad event over which no uniform-convergence statement can even be phrased.  The softmax is precisely the smoothing that restores measurability, and so the choice between soft and hard attention is, at the level of \emph{objects}, the choice of whether the learner's hypotheses form a class to which the theory of this book applies at all.
\end{remark}

\section{Exercises and Open Problems}
\label{sec:ch1-open}

The following thirteen problems are research-grade, each anchored to a concept introduced above and tagged by its epistemic character: a \emph{known unknown} (\textsc{ku}) is an articulated question awaiting an answer; an \emph{unknown known} (\textsc{uk}) is a regularity in the development or literature not yet absorbed into a clean statement.

\begin{openprob}[Sample compression conjecture, \textsc{ku}]
\label{op:compression-objects}
\Cref{thm:compression-char} characterizes finite VC dimension by compression \emph{with} side information. Prove or refute the Littlestone--Warmuth conjecture: every class of VC dimension $d$ admits a sample compression scheme of size $O(d)$ \emph{without} side information. The gap is isolated in the verified development as \leanname{FLT_Proofs.Complexity.Structures.Open_NoInfoCompressionStrengthening}; the best known unconditional bound is exponential in $d$, and $O(d)$ would tighten every compression-based generalization bound in \Cref{ch:compression}.
\end{openprob}

\begin{openprob}[The optimal No-Free-Lunch constant, \textsc{ku}]
\label{op:nfl-constant}
\Cref{lem:nfl-engine} forces disagreement on more than a quarter of a shattered set, and \cref{thm:nfl} yields expected error $\geq 1/4$. Determine the exact minimax expected error as a function of the ratio $m/|T|$ of sample size to shattered-set size, and identify the largest constant replacing $1/4$ in the averaged bound. The verified engine proves the factor-$4$ form; is $1/2 - o(1)$ attainable as $m/|T| \to 0$?
\end{openprob}

\begin{openprob}[Which finite-VC classes have finite star number? \textsc{ku}]
\label{op:star}
\Cref{prop:star-ge-vc} gives $\VC \leq \mathrm{star}$, and \cref{rmk:star-not-equiv} shows the gap can be infinite (singletons: $\VC = 1$, $\mathrm{star} = \infty$). Give a structural characterization (ideally a forbidden-subclass criterion) of the concept classes with finite VC dimension and finite star number. Because the star number controls active-learning label complexity~\cite{HannekeYang2015}, such a characterization would delimit exactly the classes for which active learning yields an exponential saving.
\end{openprob}

\begin{openprob}[An inductive-bias dimension, \textsc{uk}]
\label{op:bias-dimension}
\Cref{cor:bias-mandatory} makes bias mandatory but offers no measure of how much bias a class encodes. Construct a complexity measure $b(\mathcal{C})$, an ``inductive-bias dimension'', that is monotone under inclusion, well-behaved under products, and lower-bounds achievable sample complexity uniformly over distributions, so that the No-Free-Lunch barrier acquires a quantitative form to match its present qualitative one. The typed bias of \cref{def:inductive-bias-typed} (a preference score with a monotonicity constraint) is the natural carrier for such a measure.
\end{openprob}

\begin{openprob}[The proper/improper gap, \textsc{ku}]
\label{op:proper-improper}
It is known that proper learning can be super-polynomially costlier than improper learning~\cite{PittValiant1988}. Exhibit a single combinatorial invariant of $\mathcal{C}$ that controls the proper-versus-improper sample (or time) complexity ratio, and characterize the classes for which the ratio is bounded. The qualitative separation is settled; the separating dimension is not.
\end{openprob}

\begin{openprob}[Model-selection equivalence witnesses, \textsc{ku}]
\label{op:model-selection}
For each edge of the triangle in \cref{fig:model-selection-triangle}, either supply the missing equivalence witness or upgrade the obstruction into a theorem of non-equivalence. Concretely: characterize the parametric model families on which two-part MML coincides with Bayesian MAP beyond the constant-Fisher-information case, and decide whether the residual discrepancy can be absorbed into the side information of a compression scheme.
\end{openprob}

\begin{openprob}[Efficient version-space representation, \textsc{ku}]
\label{op:version-space}
\Cref{prop:version-space} shows the version space shrinks monotonically and never increases Littlestone dimension, yet for most classes it cannot be efficiently represented or sampled. Characterize the classes whose version space admits a polynomial-size representation under arbitrary consistent data, and relate this property to the eluder dimension~\cite{RussoVanRoy2013} and the star number (\cref{def:star-number}).
\end{openprob}

\begin{openprob}[The coding-theorem constant, \textsc{uk}]
\label{op:solomonoff}
The coding theorem $K(x) = -\log M(x) + O(1)$ (\cref{def:alg-prob}) hides an additive constant depending on the universal machine. Formalize the theorem in a proof assistant with the constant made explicit as a function of the reference machine's description, closing the gap between the algorithmic-probability and Kolmogorov-complexity foundations of MDL. No verified companion currently exists for this result.
\end{openprob}

\begin{openprob}[Constructive necessity in the fundamental theorem, \textsc{uk}]
\label{op:asymmetry}
\Cref{thm:vc-char} is proved with a constructive sufficiency direction (empirical risk minimization) and a non-constructive necessity direction (\cref{rmk:vc-asymmetry}). Give a constructive necessity proof: an explicit algorithm that, from oracle access to a class of infinite VC dimension, produces the hard distribution witnessing non-learnability with a stated dependence on the sample budget. This would remove the asymmetry noted at \leanname{FLT_Proofs.Theorem.PAC.vc_characterization} and upgrade the existential lower bounds to effective ones.
\end{openprob}

\begin{openprob}[A functorial account of the thirteen measures, \textsc{ku}]
\label{op:dimension-multiplicity}
\Cref{fig:concept-class-hub} shows 13 complexity measures targeting the concept class. Determine whether they organize into a single partial order, or a functor from the category of concept classes (with label-preserving maps) to $\N \cup \{\infty\}$, and classify the pairs that admit two-sided polynomial comparisons against the pairs that are provably incomparable. A clean answer would explain why no measure subsumes the others while many pairs are nonetheless linked.
\end{openprob}

\begin{openprob}[The VC dimension of attention, \textsc{ku}]
\label{op:attention-vc-objects}
\Cref{prop:attention-wellbehaved} shows the softmax-attention class $\mathcal{A}_{d,n_K}$ is a well-behaved object, so \cref{thm:vc-char} applies once its VC dimension is known to be finite.  Is $\VC(\mathcal{A}_{d,n_K})$ finite for every $d, n_K$, and what is its exact growth in $d$ and $n_K$?  A Pfaffian-function argument suggests finiteness, but no machine-checked bound exists; closing it would place attention definitively on the learnable side of the hub diagram (\cref{fig:concept-class-hub}) and is the verified counterpart to the well-behavedness handle \leanname{softAttention_wellBehaved}.
\end{openprob}

\begin{openprob}[A legitimate object for hard attention, \textsc{uk}]
\label{op:hard-attention-object}
\Cref{prop:soft-hard-borel} shows that argmax (hard) attention can fall outside the measurable concept classes, with an analytic-not-Borel bad event (\leanname{soft_vs_hard_attention_measurabilitySeparation}).  Identify a refined class of objects, an o-minimal or otherwise tame family of hypothesis classes, that contains hard attention and to which a uniform-convergence theory still applies, so that hard attention re-enters the theory as a legitimate learning object and sheds its status as a measurability pathology.
\end{openprob}

\begin{openprob}[Attention among the thirteen measures, \textsc{ku}]
\label{op:attention-hub}
\Cref{fig:concept-class-hub} measures the concept class with thirteen instruments.  Compute, for the softmax-attention class, the values of as many of them as possible (VC dimension, Littlestone dimension, Rademacher complexity, compression size) and determine which instruments agree and which diverge on this single modern object.  The well-behavedness of \cref{prop:attention-wellbehaved} guarantees the measures are defined; whether attention is a place where they coincide or a place where they pull apart is open.
\end{openprob}

\subsection*{Counterfactual exercises}

Each exercise perturbs the wheel of \Cref{fig:concept-class-hub} or the compression spine of \Cref{fig:model-selection-triangle} and asks where the perturbed claim lands. State for each whether the listed conditions still coincide, and name the result that settles it.

\begin{enumerate}[leftmargin=*,itemsep=3pt]
\item \textbf{Infinite label space.} Let labels range over many classes, so the hub measure becomes the DS dimension and the criterion multiclass PAC learnability. Show that finiteness of the Natarajan dimension no longer suffices, and that the characterizing quantity is the DS dimension: a class is multiclass PAC learnable iff its DS dimension is finite (Brukhim et al., 2022).
\item \textbf{Online learnability.} This is the nested inner wheel (\leanname{littlestone_characterization}), joined to the outer one by a one-way bridge \leanname{online_strictly_stronger_pac}: the online demand implies the statistical one and not conversely, with threshold rules standing at the gap.
\item \textbf{Label noise.} Drop realizability and measure success against the best rule in the class. \leanname{AgnosticPACLearnable} is defined but carries no equivalence theorem, so the classical agnostic characterization (Vapnik 1998) enters by citation.
\item \textbf{Partial concepts.} Allow partially-defined concepts. Here the central edge \emph{breaks}: there are partial classes learnable with infinite VC dimension (Alon--Hanneke--Holzman--Moran 2021), so the VC-to-learnability face fails outright.
\item \textbf{Drop measurability.} Remove the well-behavedness side-condition. A hard argmax-routed attention head has a failure event that escapes measurability (\leanname{soft_vs_hard_attention_measurabilitySeparation}), so uniform convergence cannot even be stated.
\item \textbf{Constructive necessity.} Ask the necessity direction to exhibit an explicit hard distribution. The sufficiency direction builds an empirical risk minimizer; the necessity direction is existential, and whether it can be made constructive is open (\cref{op:asymmetry}).
\item \textbf{Non-computable prior.} Replace the prior by one that is not lower-semicomputable. This breaks the Bayesian approximation edge (\leanname{bayesianToBatch}) while leaving the compression spine intact, since the spine never mentions a prior.
\item \textbf{Compression without side information.} Demand a summary of size linear in the dimension with no tag. The reverse direction is verified (\leanname{compression_bounds_vcdim}); the forward strengthening is the open Littlestone--Warmuth conjecture (\cref{op:compression-objects}).
\item \textbf{PAC-Bayes bridge.} \leanname{pac_bayes_finite} bounds the Gibbs risk by empirical risk plus a prior-weighted complexity term, so the quantity the three model-selection principles share is itself a generalization theorem.
\end{enumerate}
